\chapter{KHẢO SÁT VÀ PHÂN TÍCH YÊU CẦU}

\section{Bối cảnh và bài toán}

Đài truyền hình, toà báo, trường đại học và bảo tàng đều đang giữ những kho video lớn: bản
tin và phóng sự đã phát sóng, tường thuật thể thao, bài giảng, phim tư liệu về hiện vật.
Phần lớn các kho này nằm rời rạc ở từng đơn vị, chỉ được gắn vài thông tin chung như tên
chương trình và ngày phát hành. Người ngoài đơn vị hầu như không biết trong kho có gì. Khi
một nhà báo cần một cảnh quay cũ cho phóng sự, một nhà nghiên cứu cần tư liệu về một sự
kiện, hay một giảng viên cần đoạn video minh hoạ cho bài giảng, họ phải liên hệ từng nơi,
chờ người giữ kho tua thủ công, rồi xin phép sử dụng bằng giấy tờ.

\textbf{Thư viện số tư liệu nghe nhìn} giải quyết bài toán này theo mô hình quen thuộc của
thư viện, nhưng cho video. Đề tài xây dựng hệ thống theo một \textit{tình huống giả định}:
\begin{itemize}
    \item Nhiều \textbf{đơn vị thành viên} cùng gửi video đã số hoá vào một kho dùng chung.
    Mỗi đơn vị đóng góp các \textbf{bộ sưu tập} và giữ bản quyền của chúng; bộ sưu tập có thể
    mở cho mọi người hoặc hạn chế (dùng lại phải có giấy phép).
    \item Video được biên mục: có nhan đề, thuộc một bộ sưu tập, được gắn các \textbf{chủ đề}
    trong một cây phân loại, và được số hoá thành khung hình, chữ trên màn hình (OCR) và lời
    nói (ASR) để tra cứu được theo nội dung.
    \item \textbf{Bạn đọc} (nhà báo, nhà nghiên cứu, sinh viên) mô tả đoạn tư liệu cần tìm
    bằng tiếng Việt, ví dụ ``tay đua áo vàng về đích, phía sau là biển hiệu Hồng Thái Phúc
    Lộc'', và nhận các khung hình khớp nhất trong vài giây.
    \item Muốn dùng lại tư liệu, bạn đọc gửi \textbf{yêu cầu khai thác} kèm mục đích sử dụng
    (phát sóng, nghiên cứu, giảng dạy, trưng bày). \textbf{Cán bộ thư viện} tìm và chọn đúng
    đoạn tư liệu, rồi \textbf{cấp giấy phép sử dụng} có thời hạn hoặc từ chối.
\end{itemize}

Trọng tâm của đề tài là \textbf{hệ quản trị dữ liệu} cho bài toán này. Cơ sở dữ liệu phải
lưu được danh mục nhiều đơn vị và hàng nghìn khung hình kèm văn bản, tìm kiếm trả lời trong
tối đa 30 giây, quản lý vòng đời yêu cầu khai thác và giấy phép một cách nhất quán, phân
quyền cho nhiều loại người dùng, ghi nhật ký mọi thao tác và đo được chất lượng tìm kiếm.
Phần lớn logic được đặt ngay trong SQL Server để vận dụng các nội dung của học phần.

\section{Các bên tham gia}

\begin{itemize}
    \item \textbf{Đơn vị thành viên}: đài truyền hình, báo, trường, bảo tàng, và chính đơn vị
    vận hành thư viện. Một đơn vị có thể vừa đóng góp bộ sưu tập, vừa có bạn đọc khai thác tư
    liệu của đơn vị khác.
    \item \textbf{Bạn đọc}: người dùng tra cứu và gửi yêu cầu khai thác. Bạn đọc thuộc một đơn
    vị hoặc là cá nhân, và được phân loại theo đối tượng (nhà báo, nhà nghiên cứu, sinh viên,
    khác) để thống kê.
    \item \textbf{Cán bộ thư viện}: xử lý yêu cầu: tìm, chọn đoạn tư liệu phù hợp, duyệt và
    cấp giấy phép hoặc từ chối.
    \item \textbf{Quản trị}: quản lý tài khoản, vai trò, theo dõi nhật ký và chất lượng tìm
    kiếm.
\end{itemize}

\section{Quy trình nghiệp vụ}

\begin{enumerate}
    \item \textbf{Đăng ký}: khách xem được trang mở đầu và tự tạo tài khoản bạn đọc; tài
    khoản cán bộ thư viện do quản trị viên cấp.
    \item \textbf{Biên mục}: đơn vị đóng góp bộ sưu tập; video được ghi nhan đề, gắn chủ đề
    và số hoá (khung hình, OCR, ASR). Trong đề tài, khâu số hoá được xem là đã có sẵn.
    \item \textbf{Tra cứu}: bạn đọc nhập mô tả; hệ thống đối chiếu với hình ảnh, chữ trên màn
    hình và lời nói, trả về các khung hình đã xếp hạng kèm bằng chứng (nguồn nào tìm thấy, xếp
    hạng mấy). Trang chi tiết khung hình cho xem ảnh lớn, khung hình lân cận, chữ OCR, lời nói
    $\pm 15$ giây và phiếu mô tả video (bộ sưu tập, đơn vị giữ bản quyền, chủ đề).
    \item \textbf{Gửi yêu cầu khai thác}: bạn đọc mô tả đoạn tư liệu cần dùng và chọn mục đích
    sử dụng. Yêu cầu mới ở trạng thái \textit{chờ xử lý}.
    \item \textbf{Chọn đoạn tư liệu}: cán bộ thư viện tìm và chọn một hay nhiều khung hình cho
    yêu cầu, có thể kèm ghi chú.
    \item \textbf{Duyệt}: cán bộ cấp giấy phép (ghi đơn vị được cấp phép, thời hạn, điều kiện)
    hoặc từ chối. Chỉ cấp phép được khi yêu cầu đã có ít nhất một đoạn tư liệu. Sau khi duyệt
    hay từ chối, yêu cầu và danh sách đoạn tư liệu bị khoá (Hình~\ref{fig:request_states}).
    \item \textbf{Theo dõi}: xuất danh sách đoạn tư liệu (video, số khung hình gốc, thời điểm)
    để trích đoạn; theo dõi giấy phép sắp hết hạn; thống kê yêu cầu theo đơn vị, mục đích.
    \item \textbf{Quản trị}: tạo tài khoản, gán vai trò, khoá tài khoản, xem nhật ký; chạy
    đánh giá chất lượng tìm kiếm.
\end{enumerate}

\begin{figure}[H]
\centering
\begin{tikzpicture}[
    state/.style={draw=blue!60!black, thick, rounded corners=6pt, fill=blue!6, align=center,
        font=\small, minimum height=1cm, minimum width=2.6cm},
    arrow/.style={-{Stealth[length=2.5mm]}, thick},
    note/.style={font=\scriptsize, align=center}]
\node[circle, fill=black, inner sep=2.5pt] (start) at (-1.6,0) {};
\node[state] (pending) at (3.2,0) {\textbf{pending}\\\footnotesize chờ xử lý};
\node[state, fill=green!8, draw=green!45!black] (approved) at (9.4,1.3) {\textbf{approved}\\\footnotesize đã cấp phép};
\node[state, fill=red!6, draw=red!60!black] (rejected) at (9.4,-1.3) {\textbf{rejected}\\\footnotesize từ chối};
\draw[arrow] (start) -- node[above, note]{\texttt{sp\_submit\_request}} (pending);
\draw[arrow] (pending) -- node[above, note, sloped]{\texttt{sp\_grant\_license}\\(cần $\geq 1$ đoạn)} (approved);
\draw[arrow] (pending) -- node[below, note, sloped]{\texttt{sp\_reject\_request}} (rejected);
\draw[arrow] (pending.south) to[out=-70, in=-110, looseness=6]
    node[below, note]{\texttt{sp\_add\_selection}, \texttt{sp\_remove\_selection}} (pending.south west);
\end{tikzpicture}
\caption{Vòng đời của một yêu cầu khai thác; trạng thái chỉ đổi qua thủ tục lưu trữ}
\label{fig:request_states}
\end{figure}

\section{Mô tả hệ thống}

Hệ thống gồm cơ sở dữ liệu SQL Server và một ứng dụng web nhỏ (Hình~\ref{fig:architecture}).
Ứng dụng chỉ là lớp giao diện mỏng: mọi thao tác dữ liệu đi qua câu lệnh SQL hoặc thủ tục
lưu trữ, và chính cơ sở dữ liệu kiểm tra lại quyền cùng trạng thái của yêu cầu.

\begin{figure}[H]
\centering
\begin{tikzpicture}[
    box/.style={draw=blue!60!black, thick, rounded corners=3pt, fill=blue!6, align=center,
        font=\small, minimum height=1cm, inner sep=5pt},
    db/.style={box, fill=orange!10, draw=orange!70!black},
    arrow/.style={-{Stealth[length=2.5mm]}, thick}]
\node[box, text width=2.3cm] (browser) {Trình duyệt\\\footnotesize giao diện 3D (Three.js) và trang HTML};
\node[box, right=1.3cm of browser, text width=4.0cm] (app) {Ứng dụng web FastAPI\\
    \footnotesize trang HTML, API JSON \texttt{/api/search},\\ \footnotesize đăng nhập, kiểm tra quyền,\\ \footnotesize hạ quyền bằng \mbox{\texttt{EXECUTE AS}}};
\node[db, right=1.3cm of app, text width=4.4cm] (sql) {\textbf{SQL Server Express}\\
    \footnotesize danh mục, kho video, chỉ mục ngược,\\
    \footnotesize yêu cầu, giấy phép, hàm, thủ tục,\\
    \footnotesize view, trigger, role, nhật ký};
\node[box, below=1.6cm of app, text width=4.0cm] (visual) {Kênh hình ảnh\\
    \footnotesize SigLIP2 (CPU) + FAISS};
\node[box, below=1.6cm of sql, text width=4.4cm, fill=gray!8, draw=gray] (files) {Video đã số hoá\\
    \footnotesize khung hình, OCR, ASR};
\draw[arrow] (browser) -- node[above, font=\scriptsize]{HTTP} node[below, font=\scriptsize]{cookie} (app);
\draw[arrow] (app) -- node[above, font=\scriptsize]{ODBC} node[below, font=\scriptsize]{EXEC} (sql);
\draw[arrow] ([xshift=-8mm]app.south) -- node[left, font=\scriptsize]{mô tả} ([xshift=-8mm]visual.north);
\draw[arrow] ([xshift=8mm]visual.north) -- node[right, font=\scriptsize, align=left]{JSON\\keyframe\_id} ([xshift=8mm]app.south);
\draw[arrow, dashed] (files) -- node[right, font=\scriptsize]{\texttt{09\_sample\_data.sql}} (sql);
\end{tikzpicture}
\caption{Kiến trúc tổng thể của hệ thống}
\label{fig:architecture}
\end{figure}

Các chức năng chính:
\begin{itemize}
    \item \textbf{Danh mục thư viện}: đơn vị thành viên, bộ sưu tập (mở hoặc hạn chế), cây chủ
    đề nhiều cấp và việc gắn chủ đề cho video; phiếu mô tả video và thống kê theo bộ sưu tập,
    theo nhánh chủ đề.
    \item \textbf{Tra cứu theo mô tả}, gộp ba nguồn bằng chứng: hình ảnh (vector SigLIP2),
    chữ trên hình (OCR) và lời nói (ASR). Việc tách từ, tra chỉ mục ngược, tính điểm và gộp
    xếp hạng đều do thủ tục \texttt{sp\_search\_footage} thực hiện. Kết quả lọc được theo chủ
    đề (gồm cả chủ đề con) và theo bộ sưu tập.
    \item \textbf{Giao diện lưu trữ 3D}: cùng một kết quả tìm kiếm được trình bày theo ba
    cách (vòng tròn, hành lang, đám mây ảnh); dạng danh sách truyền thống vẫn được giữ.
    \item \textbf{Yêu cầu khai thác và giấy phép}: gửi yêu cầu, chọn đoạn tư liệu, cấp phép
    hoặc từ chối, theo dõi hiệu lực giấy phép; bạn đọc chỉ thấy yêu cầu của mình.
    \item \textbf{Tài khoản và phân quyền}: khách tự đăng ký tài khoản bạn đọc; quản trị viên
    tạo tài khoản (kèm đơn vị, đối tượng), gán vai trò, khoá tài khoản, đặt lại mật khẩu; mọi
    thay đổi được ghi nhật ký.
    \item \textbf{Đo chất lượng tìm kiếm}: chạy toàn bộ yêu cầu mẫu đã xác minh, lưu kết quả
    và tính Recall@K, MRR, độ trễ bằng view.
\end{itemize}

\section{Người dùng và vai trò}

Hệ thống phân quyền theo vai trò (RBAC). Quyền được gán cho vai trò, vai trò được gán cho
người dùng; một người có thể có nhiều vai trò (Bảng~\ref{tab:roles}). Vai trò quyết định
người dùng \textit{được làm gì}; còn đối tượng bạn đọc (nhà báo, nhà nghiên cứu, sinh viên)
và đơn vị công tác là thông tin thống kê và được ghi vào giấy phép, không cấp thêm quyền.
Tài khoản tự đăng ký luôn chỉ có vai trò \texttt{reader}; muốn trở thành cán bộ thư viện phải
được quản trị viên gán vai trò.

\begin{table}[H]
\centering
\caption{Vai trò và quyền nghiệp vụ}
\label{tab:roles}
\small
\setlength{\tabcolsep}{4pt}
\begin{tabular}{lcccccc}
\toprule
\textbf{Vai trò} & \rotatebox{70}{\texttt{\footnotesize search.run}} & \rotatebox{70}{\texttt{\footnotesize request.submit}}
 & \rotatebox{70}{\texttt{\footnotesize selection.write}} & \rotatebox{70}{\texttt{\footnotesize license.grant}}
 & \rotatebox{70}{\texttt{\footnotesize eval.view}} & \rotatebox{70}{\texttt{\footnotesize user.manage}} \\
\midrule
reader -- bạn đọc & \checkmark & \checkmark & & & & \\
librarian -- cán bộ thư viện & \checkmark & \checkmark & \checkmark & \checkmark & & \\
admin -- quản trị & \checkmark & \checkmark & \checkmark & \checkmark & \checkmark & \checkmark \\
\bottomrule
\end{tabular}
\end{table}

\section{Dữ liệu}

Đề tài không có kho video của các đơn vị thật. Phần video dùng một bộ video tiếng Việt công
khai đã được trích khung hình, OCR và ASR; toàn bộ bộ dữ liệu có 873 video với 401.153 khung
hình, 269.634 dòng OCR và 74.049 đoạn lời nói. Để giảm dung lượng phải tải, cơ sở dữ liệu mẫu
chỉ dùng \textbf{lô L23} (khoảng 1,5\% bộ dữ liệu): 25 đoạn tường thuật một giải đua xe đạp
truyền hình, có nhiều chữ trên màn hình và lời bình luận (Bảng~\ref{tab:source_data}). Cả 25
video được xếp vào một bộ sưu tập ``Tường thuật đua xe đạp truyền hình'' do đơn vị vận hành
thư viện quản lý.

Nhóm tự viết 120 \textbf{yêu cầu khai thác mẫu} theo cách một bạn đọc mô tả đoạn tư liệu cần
dùng, rồi xác minh thủ công vị trí đúng (video và thời điểm) của từng yêu cầu. Bộ yêu cầu này
vừa là dữ liệu mẫu cho quy trình khai thác, vừa dùng để đo chất lượng tìm kiếm.

\begin{table}[H]
\centering
\caption{Dữ liệu nguồn và phần được đưa vào cơ sở dữ liệu}
\label{tab:source_data}
\renewcommand{\arraystretch}{1.2}
\small
\begin{tabularx}{\textwidth}{>{\raggedright\arraybackslash}p{4.0cm} >{\raggedright\arraybackslash}X r}
\toprule
\textbf{Nguồn} & \textbf{Nội dung} & \textbf{Số dòng (L23)} \\
\midrule
\texttt{frames.csv} & Danh mục khung hình: video, số thứ tự, thời điểm, chỉ số khung hình gốc & 5.913 \\
\texttt{map-keyframes} & Tốc độ khung hình (fps) của từng video & 25 video \\
\texttt{ocr.sqlite} & Chữ nhận dạng trên khung hình & 3.976 \\
\texttt{asr.sqlite} & Lời nói theo đoạn thời gian & 1.238 \\
\texttt{l23\_requests.csv} & Yêu cầu khai thác mẫu kèm vị trí đã xác minh & 120 \\
Ảnh \texttt{L23.zip} & Ảnh khung hình để hiển thị (1,1 GB) & 5.913 tệp \\
Chỉ mục SigLIP2 & Vector 1152 chiều của khung hình (chỉ phần L23) & 5.913 \\
\bottomrule
\end{tabularx}
\end{table}

Phần danh mục bao quanh kho video là \textbf{dữ liệu minh hoạ}, được ghi rõ như vậy trong
script \texttt{09\_sample\_data.sql}:
\begin{itemize}
    \item Bốn đơn vị thành viên khác (``Đài truyền hình mẫu A'', ``Báo điện tử mẫu B'',
    ``Trường đại học mẫu C'', ``Bảo tàng mẫu D''), mỗi đơn vị có một bộ sưu tập chưa nạp
    video. Tên đều chứa chữ ``mẫu'' và không mô tả tổ chức có thật nào.
    \item Năm tài khoản demo (quản trị, một cán bộ thư viện, ba bạn đọc thuộc hai đơn vị),
    mục đích, trạng thái của 120 yêu cầu và 30 giấy phép. Ngày gửi yêu cầu và ngày cấp phép
    được tính lùi từ thời điểm nạp dữ liệu, nên có cả giấy phép còn hạn, sắp hết hạn và đã
    hết hạn.
    \item \textbf{Chủ đề} của video thì không bịa: chúng được gắn bằng luật trên lời nói thật.
    Ví dụ, video thuộc chủ đề ``Nước rút'' nếu có ít nhất một đoạn lời bình chứa cụm ``nước
    rút'' (sau khi bỏ dấu). Mọi nhãn chủ đề vì thế truy được về câu nói đã sinh ra nó.
\end{itemize}

\section{Yêu cầu phi chức năng}

\begin{itemize}
    \item \textbf{Thời gian phản hồi}: mỗi lượt tìm xong trong 30 giây; quá hạn thì câu lệnh
    bị huỷ trên máy chủ và người dùng nhận thông báo thay vì chờ mãi.
    \item \textbf{Toàn vẹn dữ liệu}: dùng đủ các loại ràng buộc PK, FK, UNIQUE, NOT NULL,
    DEFAULT, CHECK, đặt tên rõ ràng; dữ liệu dẫn xuất được tính bằng cột tính toán hoặc view;
    trạng thái yêu cầu và giấy phép luôn khớp nhau.
    \item \textbf{Bảo mật và truy vết}: mật khẩu băm có muối, phiên đăng nhập chỉ lưu giá trị
    băm; ứng dụng làm việc với quyền tối thiểu; quyền nghiệp vụ và trạng thái yêu cầu được
    kiểm tra lại trong CSDL; bạn đọc chỉ thấy yêu cầu của mình; mọi thao tác chọn đoạn tư
    liệu, duyệt yêu cầu và thay đổi tài khoản đều có nhật ký.
    \item \textbf{Chạy trên SQL Server Express}: không dùng tính năng chỉ có ở bản trả phí;
    script chạy được trên SQL Server 2017 trở lên.
    \item \textbf{Tái lập}: một tệp \texttt{full\_setup.sql} dựng lại toàn bộ CSDL và dữ liệu
    mẫu; mọi số liệu trong báo cáo đều chạy lại được.
\end{itemize}

\section{Lựa chọn công nghệ và phạm vi}

\textbf{SQL Server Express} là hệ quản trị CSDL theo yêu cầu đề bài. Bản Express không có
dịch vụ Full-Text Search, nên đề tài tự xây \textbf{chỉ mục ngược} bằng bảng: từ điển từ
khoá và danh sách (từ khoá, khung hình). Đây cũng là cấu trúc cốt lõi của các hệ thống tìm
kiếm văn bản~\cite{manning2008ir}, nên thể hiện rõ vai trò của chỉ mục và khoá trong CSDL.

Kênh hình ảnh dùng mô hình đa ngôn ngữ \textbf{SigLIP2}~\cite{tschannen2025siglip2} (kế thừa
SigLIP~\cite{zhai2023siglip}) để biến mô tả tiếng Việt thành vector, rồi tra chỉ mục
\textbf{FAISS}~\cite{johnson2021faiss} chứa sẵn vector của các khung hình. Phần này chạy ngoài
CSDL vì SQL Server không có kiểu vector; kết quả (danh sách \texttt{keyframe\_id} đã xếp hạng)
được đưa vào thủ tục tìm kiếm dưới dạng JSON để gộp với các nguồn văn bản. Các danh sách xếp
hạng được gộp bằng \textbf{Reciprocal Rank Fusion}~\cite{cormack2009rrf}.

Ứng dụng web viết bằng Python (FastAPI), kết nối SQL Server qua ODBC. Phạm vi đề tài không
gồm khâu số hoá video (trích khung hình, OCR, ASR), vốn được xem là đã có sẵn; không gồm việc
thu phí, ký hợp đồng bản quyền với đơn vị đóng góp, hay chuyển tệp video gốc cho bạn đọc sau
khi được cấp phép.
